\section{Problem Formulation and Fixed Benchmark}
\label{sec:problem}

\subsection{Hamiltonian and fragmentation}

Let $\Lambda=(\mathbb{Z}/L\mathbb{Z})^2$ be a periodic square lattice with one
qubit per site.  The target Hamiltonian is
\begin{equation}
 H=\sum_{\langle u,v\rangle}h_{uv},\qquad
 h_{uv}=\frac{X_uX_v+Y_uY_v+Z_uZ_v}{4}.
 \label{eq:hamiltonian}
\end{equation}
Horizontal and vertical bonds are split by parity into four perfect
matchings $M_1,\ldots,M_4$.  Defining
$H_a=\sum_{(u,v)\in M_a}h_{uv}$ gives
$H=H_1+H_2+H_3+H_4$.  Terms within each matching have disjoint support and
therefore commute exactly.  For real-time analysis we use anti-Hermitian
generators $A_a=-iH_a$ and $A=\sum_a A_a$.

The principal benchmark fixes
\begin{equation}
 L=12,qquad N=L^2=144,qquad T=1,qquad \eps=10^{-6}.
 \label{eq:benchmark}
\end{equation}
The required guarantee is uniform over all input states:
\begin{equation}
 \opnorm{e^{-iTH}-\widetilde U_r(T)}\le \eps,
 \label{eq:uniformerror}
\end{equation}
where $\opnorm{\cdot}$ is the spectral norm.  This is deliberately stronger
than an observable-specific, state-dependent, or average-case condition.

\subsection{Product formula}

Let $S_2(t)$ denote the symmetric second-order formula for the four fragments.
The standard five-copy Suzuki recursion \citep{suzuki1991general} forms
\begin{equation}
 S_4(t)=S_2(pt)^2S_2((1-4p)t)S_2(pt)^2,\qquad
 p=(4-4^{1/3})^{-1}.
 \label{eq:suzuki4}
\end{equation}
After merging adjacent equal fragments, one step contains 31 group
exponentials.  The $r$-step approximation is
\begin{equation}
 \widetilde U_r(T)=S_4(T/r)^r.
\end{equation}
The baseline and candidate use precisely the same ordered stage sequence and
the same outward interval for the algebraic coefficient $p$.  Consequently,
any reduction in $r$ is a proof improvement rather than a change of
simulation formula.

\subsection{Cost model}

Equal fragments at the boundary between adjacent symmetric steps merge.  The
first step contributes 31 groups and each additional step contributes 30:
\begin{equation}
 G(r)=31+30(r-1)=30r+1.
 \label{eq:groups}
\end{equation}
Each perfect matching contains $N/2=72$ disjoint bonds.  We therefore count
\begin{equation}
 B(r)=72G(r)
 \label{eq:bonds}
\end{equation}
bond propagators.  Using the same conservative upper bound of three CNOTs per
normalized Heisenberg bond propagator on both sides gives
\begin{equation}
 C(r)\le 3B(r)=216G(r).
 \label{eq:cnots}
\end{equation}
The primary ratio is identical for groups, bond propagators, and this CNOT
upper bound because the multiplicative factors cancel.

\subsection{Compiler contract}

The compiler receives a tuple
\begin{equation}
 \mathcal I=(H,\mathcal F,S,T,\eps,\mathcal C),
\end{equation}
where $\mathcal F$ is a fragmentation, $S$ a product formula, and
$\mathcal C$ a cost model.  It must return a tuple
\begin{equation}
 \mathcal O=(r_\star,E_{r_\star},E_{r_\star-1},R_\star,\cert),
\end{equation}
such that an independent verifier establishes
\begin{align}
 E_{r_\star}&\le\eps,& E_{r_\star-1}&>\eps,\label{eq:adjacent}\
 R_\star&=\mathcal C(S,r_\star),
\end{align}
using exact integers, rational numbers, or outward rational intervals.  The
second inequality does not prove that the physical circuit with
$r_\star-1$ fails; it proves that $r_\star$ is the smallest integer accepted
by this particular certified upper-bound function.  This distinction is
retained throughout.
